\section{场内容与手征规范结构}\label{sec:sm-fields-gauge-dp11}

QED 中一个 Dirac 电子的左、右手分量具有相同电荷；QCD 中同一个夸克又同时携带颜色指标。这已经说明一个场可以同时属于若干彼此独立的内部表示。标准模型只把这个事实再推进一步：弱 $SU(2)_L$ 不再同样作用于两个手征分量，因此在写协变导数以前，必须先明确“每一个 Weyl 场究竟属于哪个内部表示”。

表示数据一旦给定，协变导数便由各群因子的生成元直接相加；裸质量项是否允许，只需检查左右手表示能否收缩成单态，Yukawa 耦合是否存在也由同一个单态条件决定。先完成这些局域表示约束以后，Higgs 真空和质量本征态才有明确的输入对象。

\subsection{费米子的手征表示}

标准模型的基本费米子都是自旋 $1/2$ 场，但弱相互作用对左手与右手分量的作用不同。这里的 Weyl 场可以理解为 Dirac 四分量场被手征投影后得到的一个不可约手征分量；它仍是量子场，而不是“某个电子朝左旋转”的经典图像。采用手征投影算符
\begin{equation*}
 P_L=\frac{1-\gamma^5}{2},
 \qquad
 P_R=\frac{1+\gamma^5}{2},
\end{equation*}
定义 $\psi_L=P_L\psi$、$\psi_R=P_R\psi$。因为 $P_LP_R=0$，一个 Dirac 场可以视为两个 Lorentz 不等价 Weyl 表示的直和；手征规范理论允许这两个手征分量属于不同的内部对称表示。

一个费米子场可以同时带颜色指标、弱同位旋指标和一个 $U(1)$ 荷。数学上，这只是直积群表示的张量积。若
\begin{equation*}
 G=G_1\times\cdots\times G_k\times U(1)_1\times\cdots\times U(1)_r,
\end{equation*}
则一个场由每个非阿贝尔因子的表示 $R_a$ 和每个阿贝尔因子的生成元本征值 $q_\rho$ 唯一指定。局域变换写成
\begin{equation*}
 \Psi(x)\mapsto
 \left[\bigotimes_aR_a(U_a(x))\right]
 \exp\!\left(\ii\sum_\rho q_\rho\theta_\rho(x)\right)\Psi(x).
\end{equation*}
这里的“表示”就是内部矩阵怎样作用在场分量上；不同群因子作用在不同张量因子上，因此彼此交换。

对标准模型，记号
\begin{equation*}
 (R_3,R_2)_Y
\end{equation*}
依次表示颜色 $SU(3)_c$ 表示、弱 $SU(2)_L$ 表示和 $U(1)_Y$ 生成元本征值。$SU(2)_L$ 的三个生成元记为 $T^I$；对二重态它们就是 $\sigma^I/2$，其中 $T^3$ 的两个本征值为 $+1/2$ 与 $-1/2$，这就是“弱同位旋第三分量”的数学含义。$Y$ 称为\emph{弱超荷}，它是 $U(1)_Y$ 变换的生成元本征值。Higgs 真空确定未破缺方向以后，物理电荷生成元才由 $Q=T^3+Y$ 给出。$T^3$、$Y$ 与最终电荷 $Q$ 因而分别对应弱同位旋、超荷与未破缺电磁方向。

对每一代费米子，定义五个 Weyl 多重态
\begin{align}
 Q_L&=\begin{pmatrix}u_L\\ d_L\end{pmatrix},
 &Q_L&:\quad(\mathbf3,\mathbf2,+\tfrac16),\\
 u_R&,&u_R&:\quad(\mathbf3,\mathbf1,+\tfrac23),\\
 d_R&,&d_R&:\quad(\mathbf3,\mathbf1,-\tfrac13),\\
 L_L&=\begin{pmatrix}\nu_L\\ e_L\end{pmatrix},
 &L_L&:\quad(\mathbf1,\mathbf2,-\tfrac12),\\
 e_R&,&e_R&:\quad(\mathbf1,\mathbf1,-1).
 \label{eq:sm-one-generation-reps-dp11}
\end{align}
\eqnrefs{eq:sm-one-generation-reps-dp11} 给出最小标准模型的一组表示数据。为了让场表示与相应规范顶角一一对应，\tabref{tab:sm-one-generation-reps-dp115} 把同一信息按“场、多重态、颜色表示、弱表示、超荷”列开。这里 $u,d,e,\nu$ 先表示一代中的上型夸克、下型夸克、带电轻子和中微子；第二、三代具有相同规范表示，只更换世代标签。

\begin{table}[htbp]
\centering
\caption{一代最小标准模型费米子的手征多重态与规范表示。电荷 $Q=T^3+Y$ 要在 Higgs 真空确定未破缺 $U(1)_{\rm em}$ 后才成为物理电荷。}
\label{tab:sm-one-generation-reps-dp115}
\begin{tabular}{c c c c c}
\toprule
场 & 多重态分量 & $SU(3)_c$ & $SU(2)_L$ & $Y$\\
\midrule
$Q_L$ & $(u_L,d_L)^T$ & $\mathbf3$ & $\mathbf2$ & $+1/6$\\
$u_R$ & $u_R$ & $\mathbf3$ & $\mathbf1$ & $+2/3$\\
$d_R$ & $d_R$ & $\mathbf3$ & $\mathbf1$ & $-1/3$\\
$L_L$ & $(\nu_L,e_L)^T$ & $\mathbf1$ & $\mathbf2$ & $-1/2$\\
$e_R$ & $e_R$ & $\mathbf1$ & $\mathbf1$ & $-1$\\
\bottomrule
\end{tabular}
\end{table}

Yukawa 相互作用要求相应场乘积成为规范单态，量子规范反常的消失又要求各手征表示的超荷满足额外代数关系；固定 $Y_\Phi=1/2$ 的归一化后，这些条件与\eqnrefs{eq:sm-one-generation-reps-dp11} 相容。最小可重整化标准模型不含 $\nu_R$，因而不存在可重整化的中微子 Dirac Yukawa 项；观测到的中微子质量必须来自额外自由度或高维算符。

\subsection{三种规范联络}

直积规范群的每一个因子都有自己的联络，而一个给定物质场只感受到它实际所属表示中的生成元。颜色联络作用在 $SU(3)_c$ 指标，弱联络作用在 $SU(2)_L$ 指标，超荷联络则乘以该场的数值超荷 $Y$；三者作用在不同内部张量因子上，所以协变导数由三项线性相加得到。明确这一点以后，任意费米子或 Higgs 场的规范顶角都可直接从其表示标签读出。

对表示 $(R_3,R_2,Y)$ 中的任意场 $\Psi$，协变导数定义为
\begin{equation}
 \boxed{
 D_\mu\Psi
 =\left[
 \partial_\mu
 +\ii g_s\mathscr C_\mu^aT_{R_3}^a
 +\ii g\mathcal W_\mu^I t_{R_2}^I
 +\ii g'Y\mathcal B_\mu
 \right]\Psi.}
 \label{eq:sm-covariant-derivative-dp11}
\end{equation}
其中 $a=1,\ldots,8$，$I=1,2,3$；对 $SU(2)$ 二重态有 $t^I=\sigma^I/2$，单态上 $t^I=0$；对颜色单态有 $T^a=0$。三个规范耦合常数
\begin{equation*}
 g_s,\qquad g,\qquad g'
\end{equation*}
分别属于 $SU(3)_c$、$SU(2)_L$ 与 $U(1)_Y$。


三个场强都由相应协变导数的对易子定义：
\begin{align}
 \mathcal F_{\mu\nu}^a
 &=\partial_\mu\mathscr C_\nu^a-\partial_\nu\mathscr C_\mu^a
 -g_sf^{abc}\mathscr C_\mu^b\mathscr C_\nu^c,\\
 \mathcal W_{\mu\nu}^I
 &=\partial_\mu\mathcal W_\nu^I-\partial_\nu\mathcal W_\mu^I
 -g\epsilon^{IJK}\mathcal W_\mu^J\mathcal W_\nu^K,\\
 \mathcal B_{\mu\nu}
 &=\partial_\mu\mathcal B_\nu-\partial_\nu\mathcal B_\mu.
 \label{eq:sm-field-strengths-dp11}
\end{align}
只有 $U(1)_Y$ 的场强不含规范场二次项；因此超荷规范玻色子在未破缺相中没有自相互作用，而弱规范玻色子与胶子具有三点和四点规范顶角。

费米子动能拉氏量对所有多重态求和：
\begin{equation}
 \Lag_{\rm fermion}
 =\sum_{\Psi=Q_L,u_R,d_R,L_L,e_R}
 \bar\Psi\,\ii\hbar c\gamma^\mu D_\mu\Psi.
 \label{eq:sm-fermion-kinetic-dp11}
\end{equation}
因为 $Q_L,L_L$ 是左手场，也可以完全采用二分量 Weyl 记号；采用四分量记号并保留 $P_L$，可使同一套 Dirac $\gamma$ 矩阵代数直接用于手征投影计算。

纯规范场部分为
\begin{equation}
 \boxed{
 \Lag_{\rm gauge}
 =-\frac{\hbar c}{4}\mathcal F_{\mu\nu}^a\mathcal F^{a\mu\nu}
 -\frac{\hbar c}{4}\mathcal W_{\mu\nu}^I\mathcal W^{I\mu\nu}
 -\frac{\hbar c}{4}\mathcal B_{\mu\nu}\mathcal B^{\mu\nu}.}
 \label{eq:sm-gauge-lagrangian-dp11}
\end{equation}
每项都具有 $\mathrm{J/m^3}$：$[\hbar c]=\mathrm{J\,m}$，$[\mathcal F]=\mathrm{m^{-2}}$。



\begin{derivation}[从\eqnrefs{eq:sm-covariant-derivative-dp11}到\eqnrefs{eq:sm-field-strengths-dp11}]
从正文\eqnrefs{eq:sm-covariant-derivative-dp11} 的协变条件出发，先求规范联络的有限变换律，再计算两个协变导数的对易子，并分别代入三个规范群因子。

把三个群因子分别写成
\begin{equation*}
 U_3(x)=\exp[\ii\alpha^a(x)T^a],\qquad
 U_2(x)=\exp[\ii\theta^I(x)t^I],\qquad
 U_Y(x)=\exp[\ii Y\vartheta_Y(x)].
\end{equation*}
对表示 $(R_3,R_2,Y)$ 中的场，局域变换为
\begin{equation*}
 \Psi'(x)=U_3(x)U_2(x)U_Y(x)\Psi(x).
\end{equation*}
三个群因子作用在颜色、弱同位旋和超荷三个不同内部空间，因而彼此交换。规范联络的变换律可先在一个一般非阿贝尔因子上推导，再分别代入 $SU(3)_c$ 与 $SU(2)_L$；阿贝尔 $U(1)_Y$ 是同一结果在对易生成元下的特例。记
\begin{equation*}
 D_\mu=\partial_\mu+\ii g_{\rm G}\mathscr C_\mu,
 \qquad \mathscr C_\mu=\mathscr C_\mu^aT^a,
\end{equation*}
并要求协变导数本身按相应表示协变变换：
\begin{equation*}
 D_\mu'\Psi'=U D_\mu\Psi.
\end{equation*}
左边逐项展开，
\begin{align*}
 D_\mu'\Psi'
 &=(\partial_\mu+\ii g_{\rm G}\mathscr C_\mu')U\Psi,\\
 &=(\partial_\mu U)\Psi+U\partial_\mu\Psi
   +\ii g_{\rm G}\mathscr C_\mu'U\Psi.
\end{align*}
右边为
\begin{equation*}
 U D_\mu\Psi
 =U\partial_\mu\Psi+\ii g_{\rm G}U\mathscr C_\mu\Psi.
\end{equation*}
消去共同的 $U\partial_\mu\Psi$ 后，任意 $\Psi$ 均须满足
\begin{equation*}
 (\partial_\mu U)+\ii g_{\rm G}\mathscr C_\mu'U
 =\ii g_{\rm G}U\mathscr C_\mu.
\end{equation*}
右乘 $U^{-1}$，得到连接的有限局域变换律
\begin{equation*}
 {
 \mathscr C_\mu'
 =U\mathscr C_\mu U^{-1}
 +\frac{\ii}{g_{\rm G}}(\partial_\mu U)U^{-1}.}
\end{equation*}
第二项是非齐次项；它正好抵消普通导数作用在局域群元上产生的 $\partial_\mu U$。若 $U$ 与 $\mathscr C_\mu$ 交换，上式退化为阿贝尔规范联络的平移。

对超荷群，直接取 $U_Y=\exp(\ii Y\vartheta_Y)$。由
\begin{align*}
 D_\mu'\Psi'
 &=\left(\partial_\mu+\ii g'Y\mathcal B_\mu'\right)
   \ee^{\ii Y\vartheta_Y}\Psi,\\
 &=\ee^{\ii Y\vartheta_Y}
 \left[\partial_\mu+\ii Y\partial_\mu\vartheta_Y
 +\ii g'Y\mathcal B_\mu'\right]\Psi
\end{align*}
与 $D_\mu'\Psi'=\ee^{\ii Y\vartheta_Y}D_\mu\Psi$ 比较，得到
\begin{equation*}
 {
 \mathcal B_\mu'=\mathcal B_\mu-\frac1{g'}\partial_\mu\vartheta_Y.}
\end{equation*}

接着由协变导数的对易子推出曲率，也就是规范场强。对任意检验场 $\Psi$，
\begin{align*}
 [D_\mu,D_\nu]\Psi
 =&\left[\partial_\mu+\ii g_{\rm G}\mathscr C_\mu,
 \partial_\nu+\ii g_{\rm G}\mathscr C_\nu\right]\Psi,\\
 =&\ii g_{\rm G}(\partial_\mu \mathscr C_\nu-\partial_\nu \mathscr C_\mu)\Psi
 -(g_{\rm G})^2[\mathscr C_\mu,\mathscr C_\nu]\Psi.
\end{align*}
利用
\begin{equation*}
 [\mathscr C_\mu,\mathscr C_\nu]
 =\mathscr C_\mu^a\mathscr C_\nu^b[T^a,T^b]
 =\ii f^{abc}\mathscr C_\mu^a\mathscr C_\nu^bT^c,
\end{equation*}
可整理为
\begin{equation*}
 [D_\mu,D_\nu]
 =\ii g_{\rm G}\mathcal F_{\mu\nu},
\end{equation*}
其中
\begin{equation*}
 {
 \mathcal F_{\mu\nu}^a
 =\partial_\mu \mathscr C_\nu^a-\partial_\nu \mathscr C_\mu^a
 -g_{\rm G}f^{abc}\mathscr C_\mu^b\mathscr C_\nu^c.}
\end{equation*}
由于上式等价于算符恒等式 $D_\mu'=UD_\mu U^{-1}$，
\begin{align*}
 [D_\mu',D_\nu']
 &=UD_\mu U^{-1}UD_\nu U^{-1}
  -UD_\nu U^{-1}UD_\mu U^{-1},\\
 &=U[D_\mu,D_\nu]U^{-1},
\end{align*}
从而
\begin{equation*}
 {\mathcal F_{\mu\nu}'=U\mathcal F_{\mu\nu}U^{-1}.}
\end{equation*}
规范联络含非齐次项，而场强按伴随表示齐次变换，因此 $\Tr(\mathcal F_{\mu\nu}\mathcal F^{\mu\nu})$ 具有规范不变性。把这一一般结果分别取 $(g_{\rm G},\mathscr C_\mu,T^a,f^{abc})=(g_s,\mathscr C_\mu,T^a,f^{abc})$ 与 $(g,\mathcal W_\mu,t^I,\epsilon^{IJK})$，便得到色规范场与弱规范场的场强；阿贝尔超荷群满足 $f^{abc}=0$，所以场强只保留普通旋度项。
\end{derivation}


\subsection{电弱表示如何约束质量项与 Yukawa 结构}

质量项是否允许首先是一个表示收缩问题。Lorentz 标量 Dirac 双线性必须连接左、右手 Weyl 分量；若这两个分量在未破缺规范群下属于不同表示，它们的乘积就还携带未收缩的内部指标。Higgs 二重态的作用正是提供这些缺失的 $SU(2)_L$ 与超荷量子数，使三场 Yukawa 算符能够收缩成规范单态。

Dirac 质量项把左右手场直接配对：
\begin{equation}
 \Lag_m=-mc^2(\bar\psi_L\psi_R+\bar\psi_R\psi_L).
 \label{eq:dirac-mass-chiral-mixing-dp11}
\end{equation}
它能够出现在未破缺规范理论中的必要条件，是 $\bar\psi_L\psi_R$ 在全部内部规范因子下都能收缩成单态。标准模型的表示表~\eqnrefs{eq:sm-one-generation-reps-dp11} 不满足这一条件：$Q_L,L_L$ 是 $SU(2)_L$ 二重态，而 $u_R,d_R,e_R$ 是单态，并且相应超荷也不同。因此未破缺的 $SU(2)_L\times U(1)_Y$ 禁止裸夸克和带电轻子 Dirac 质量项。

因此需要一个额外标量表示来补齐 $\bar\psi_L(\cdots)\psi_R$ 缺少的弱同位旋与超荷指标，使整个三场乘积在未破缺电弱群下收缩成单态。满足这一要求的最小选择是复标量二重态
\begin{equation}
 \boxed{
 \Phi=\begin{pmatrix}\phi^+\\ \phi^0\end{pmatrix},
 \qquad
 \Phi:(\mathbf1,\mathbf2)_{+1/2}.}
 \label{eq:sm-higgs-rep-dp11}
\end{equation}
它不带颜色，带一个弱二重态指标和超荷 $+1/2$。由第一章已经建立的 $SU(2)$ 伪实结构，可定义共轭二重态
\begin{equation*}
 \widetilde\Phi\equiv\ii\sigma^2\Phi^*,
 \qquad
 \widetilde\Phi:(\mathbf1,\mathbf2)_{-1/2}.
\end{equation*}
因此 $\Phi$ 与 $\widetilde\Phi$ 分别能够补偿下型和上型费米子所需的超荷差，同时都提供一个弱二重态指标。

逐项检查 $SU(2)_L$ 指标与超荷守恒后，最低维的规范不变费米子--标量组合只能取
\begin{equation}
 \boxed{
 \bar L_L\Phi e_R,
 \qquad
 \bar Q_L\Phi d_R,
 \qquad
 \bar Q_L\widetilde\Phi u_R,}
 \label{eq:sm-yukawa-invariant-structures-dp43}
\end{equation}
以及它们的厄米共轭。\eqnrefs{eq:sm-yukawa-invariant-structures-dp43} 是 Yukawa 拉氏量的表示论来源；Higgs 场取得非零真空期望值后，这些三场相互作用才转化为费米子质量矩阵。最小可重整化场内容没有 $\nu_R$，所以同一机制不能生成中微子 Dirac 质量；中微子质量需要额外场或高维算符。


\begin{figure}[htbp]
\centering
\begin{tikzpicture}[node distance=1.55cm and 2.0cm,>=Latex,line label/.style={fill=white,fill opacity=.96,text opacity=1,inner sep=1.0pt}]
\node[draw,rounded corners,align=center] (rep) {chiral representations\\$Q_L,L_L,u_R,d_R,e_R$};
\node[draw,rounded corners,right=of rep,align=center] (mass) {bare Dirac mass\\not a gauge singlet};
\node[draw,rounded corners,below=of mass,align=center] (h) {add $\Phi:(\mathbf1,\mathbf2)_{1/2}$};
\node[draw,rounded corners,left=of h,align=center] (y) {gauge-invariant\\Yukawa operators};
\draw[->] (rep)--node[line label,midway,above,font=\scriptsize]{compare $L/R$ reps}(mass);
\draw[->] (mass)--node[line label,midway,right,font=\scriptsize]{missing index}(h);
\draw[->] (h)--node[line label,midway,below,font=\scriptsize]{contract reps}(y);
\draw[->] (rep)--node[line label,midway,left,font=\scriptsize]{same reps}(y);
\end{tikzpicture}
\caption{Higgs 表示的表示论来源。手征费米子的左右表示使裸 Dirac 质量无法收缩成规范单态；引入能够补齐弱同位旋与超荷指标的最低维标量二重态后，规范单态条件确定相应 Yukawa 算符。}
\label{fig:sm-representation-to-yukawa-dp87}
\end{figure}
\figref{fig:sm-representation-to-yukawa-dp87} 应从左上角开始读：先比较左右手费米子的表示，裸质量项因此不是规范单态；Higgs 二重态提供缺失的弱同位旋与超荷指标，Yukawa 三场算符随后才能写成规范单态。

为了判断还有哪些局域项能够进入四维基本拉氏量，先把 SI 拉氏量密度换成按逆长度计数的形式，
\begin{equation}
 \widehat{\Lag}\equiv\frac{\Lag}{\hbar c},
 \qquad
 S=\hbar\int\dd^4 x\,\widehat{\Lag},
 \qquad \Delta_\partial=1.
 \label{eq:sm-geometric-lagrangian-dp68}
\end{equation}
由规范化动能项确定的工程维数为
\begin{equation}
 \boxed{
 \Delta_\Phi=1,\qquad
 \Delta_\psi=\frac32,\qquad
 \Delta_{\rm conn}=1,\qquad
 \Delta_{\mathfrak F}=2.}
 \label{eq:sm-field-engineering-dimensions-dp68}
\end{equation}
因此四维微扰可重整化局域算符必须同时满足 Lorentz 标量、$G_{\rm SM}$ 单态和 $\Delta_{\mathcal O}\le4$。结合上述场维数，独立的算符类别可概括为
\begin{equation}
 \boxed{
 \mathcal O_{\Delta\le4}\in
 \left\{
 \mathfrak F_{\mu\nu}\mathfrak F^{\mu\nu},
 \bar\psi\ii\gamma^\mu D_\mu\psi,
 (D_\mu\Phi)^\dagger D^\mu\Phi,
 \Phi^\dagger\Phi,
 (\Phi^\dagger\Phi)^2,
 \bar\psi_L\Phi\psi_R,
 \bar\psi_L\widetilde\Phi\psi_R
 \right\},}
 \label{eq:sm-renormalizable-operator-classes-dp68}
\end{equation}
这里 $\Delta_{\rm conn}$ 表示任意几何归一化规范联络的工程维数；算符分类中的抽象规范曲率统一记为 $\mathfrak F_{\mu\nu}$，而 QED 章节的物理 SI 电磁场强始终记为 $F_{\mu\nu}$。费米子双线性只有在左右手内部表示能够收缩成规范单态时才允许；标准模型表示恰好排除了裸 Dirac 质量而保留\eqnrefs{eq:sm-yukawa-invariant-structures-dp43} 的 Yukawa 配对。

\begin{derivation}[从\eqnrefs{eq:dirac-mass-chiral-mixing-dp11}到\eqnrefs{eq:sm-yukawa-invariant-structures-dp43}]
以带电轻子为例，$L_L$ 与 $e_R$ 的表示分别为

\begin{equation*}
 L_L:(\mathbf1,\mathbf2,-\tfrac12),
 \qquad
 e_R:(\mathbf1,\mathbf1,-1).
\end{equation*}
共轭场 $\bar L_L$ 属于 $SU(2)$ 共轭二重态，超荷符号反转。因此裸双线性满足
\begin{equation*}
 \bar L_L e_R:
 \qquad
 \bar{\mathbf2}\otimes\mathbf1=\bar{\mathbf2},
 \qquad
 Y=+\frac12-1=-\frac12.
\end{equation*}
它既不是 $SU(2)_L$ 单态，也不是 $U(1)_Y$ 中性量，因而不能出现在未破缺电弱理论的局域拉氏量中。夸克左右手场具有同样的二重态--单态错配，所以裸夸克 Dirac 质量也被排除。

插入 Higgs 二重态后，带电轻子组合的弱同位旋部分为
\begin{equation*}
 \bar{\mathbf2}\otimes\mathbf2
 =\mathbf1\oplus\mathbf3,
\end{equation*}
其中包含单态；取指标收缩 $\bar L_{Li}\Phi_i$ 即选出这一单态。其总超荷为
\begin{equation*}
 -Y(L_L)+Y(\Phi)+Y(e_R)
 =+\frac12+\frac12-1=0,
\end{equation*}
故 $\bar L_L\Phi e_R$ 在全部电弱规范因子下不变。

对下型夸克，弱同位旋指标由 $\bar Q_{Li}\Phi_i$ 收缩成 $SU(2)_L$ 单态。其总超荷为
\begin{align*}
 Y(\bar Q_L\Phi d_R)
 &=-Y(Q_L)+Y(\Phi)+Y(d_R),\\
 &=-\frac16+\frac12-\frac13,\\
 &=0.
\end{align*}
颜色指标则由 $\bar{\mathbf3}\otimes\mathbf3$ 中的单态收缩，因此 $\bar Q_L\Phi d_R$ 在三个规范因子下都为单态。上型夸克若使用 $\Phi$，超荷不能抵消；改用
\begin{equation*}
 \widetilde\Phi:(\mathbf1,\mathbf2,-\tfrac12)
\end{equation*}
后，
\begin{equation*}
 -Y(Q_L)+Y(\widetilde\Phi)+Y(u_R)
 =-\frac16-\frac12+\frac23=0,
\end{equation*}
于是 $\bar Q_L\widetilde\Phi u_R$ 也是规范单态。这三种组合恰好给出\eqnrefs{eq:sm-yukawa-invariant-structures-dp43}。
\end{derivation}


